\documentclass[11pt]{article}

\usepackage[a4paper,margin=1in]{geometry}
\usepackage{amsmath,amssymb,amsthm,mathtools}
\usepackage{microtype}
\usepackage[colorlinks=true,linkcolor=blue,citecolor=blue,urlcolor=blue]{hyperref}

\newtheorem{theorem}{Theorem}[section]
\newtheorem{proposition}[theorem]{Proposition}
\newtheorem{lemma}[theorem]{Lemma}
\newtheorem{corollary}[theorem]{Corollary}
\newtheorem{remark}[theorem]{Remark}

\newcommand{\R}{\mathbb{R}}
\newcommand{\one}{\mathbf{1}}
\newcommand{\ip}[2]{\left\langle #1,#2\right\rangle}
\newcommand{\Tr}{\operatorname{Tr}}
\newcommand{\op}{\mathrm{op}}
\newcommand{\HS}{\mathrm{HS}}
\newcommand{\Alt}{\mathrm{alt}}

\title{Semi-inducibility of Alternating Cycles with length \texorpdfstring{$4k+2$}{4k+2}}
\author{Taeyoung Kim}
\date{\today}

\begin{document}
\maketitle

\begin{abstract}
Let $W\colon[0,1]^2\to[0,1]$ be a graphon, and let $C_{2m}^{\Alt}$ denote the $2m$-cycle in which red edges contribute $W$ and blue edges contribute its complement $1-W$.  For every odd integer $m\ge3$, we first prove the universal inequality
\[
   4^m t(C_{2m}^{\Alt},W)+t(C_{2m},2W-1)\le1,
\]
and hence $t(C_{2m}^{\Alt},W)\le4^{-m}$.  We then refine this result at fixed edge density.  Suppose
$W$ has edge density $p$, and let $q=1-p$.  If
\[
   \frac{5-\sqrt5}{10}\le p\le \frac{5+\sqrt5}{10},
\]
then
\[
   t(C_{2m}^{\Alt},W)+t(C_{2m},W-p)\le \bigl(p(1-p)\bigr)^m.
\]
Since $t(C_{2m},W-p)\ge0$, this yields the exact fixed-edge-density profile
\[
   \sup_{t(K_2,W)=p}t(C_{2m}^{\Alt},W)=\bigl(p(1-p)\bigr)^m,
\]
with equality only for the constant graphon $W\equiv p$.  Thus the universal result holds at every edge density, while the exact fixed-density profile is determined throughout the central interval.  The proof uses normalized kernel operators, a rank-two determinant factorization, monotone auxiliary coefficient sequences, and an odd logarithmic coefficient inequality.
\end{abstract}

\section{Introduction and main result}

Alternating cycles arise in the semi-inducibility problem studied by Basit, Granet, Horsley,
K\"undgen, and Staden~\cite{BGHKS}.  They obtained sharp or nearly sharp results for alternating
cycles whose lengths are divisible by four, while the lengths congruent to two modulo four remained
open.  The first such case, the alternating six-cycle, was settled by Chen and
Noel~\cite{ChenNoel}.  Theorem~\ref{thm:universal} resolves the entire congruence class without an
edge-density restriction.  We then impose a fixed edge density and determine the sharper exact
profile throughout the central interval in Theorem~\ref{thm:main}.  For background on graphons and
homomorphism densities, see Lov\'asz~\cite{Lovasz}.

A \emph{graphon} is a symmetric measurable function
\[
   W\colon[0,1]^2\to[0,1].
\]
Its edge density is
\[
   p=t(K_2,W)=\int_{[0,1]^2}W(x,y)\,dx\,dy,
\]
and throughout we write
\[
   q:=1-p.
\]
For any bounded symmetric kernel $K$, let $T_K$ denote the associated integral operator on $L^2([0,1])$,
\[
   (T_K f)(x)=\int_0^1 K(x,y)f(y)\,dy.
\]

Fix an odd integer $m\ge3$.  The alternating red--blue cycle density is
\begin{align}
 t(C_{2m}^{\Alt},W)
 :=\int_{[0,1]^{2m}}
 \prod_{j=1}^{m}
   W(x_{2j-1},x_{2j})
   \bigl(1-W(x_{2j},x_{2j+1})\bigr)
 \,dx_1\cdots dx_{2m},
 \label{eq:alt-density}
\end{align}
where $x_{2m+1}=x_1$.
For a bounded symmetric signed kernel $K$, we use the usual cycle-density notation
\[
   t(C_r,K)
   =\int_{[0,1]^r}\prod_{j=1}^r K(x_j,x_{j+1})\,dx_1\cdots dx_r,
   \qquad x_{r+1}=x_1.
\]

We first state the density-independent result.

\begin{theorem}[Universal alternating-cycle bound]\label{thm:universal}
Let $m\ge3$ be odd.  Every graphon $W$ satisfies
\begin{equation}
   4^m t(C_{2m}^{\Alt},W)+t(C_{2m},2W-1)\le1.
   \label{eq:universal-strong}
\end{equation}
Consequently,
\begin{equation}
   t(C_{2m}^{\Alt},W)\le4^{-m}.
   \label{eq:universal-bound}
\end{equation}
Equality in \eqref{eq:universal-bound} holds if and only if $W=1/2$ almost everywhere.
\end{theorem}

At fixed edge density, the universal upper bound can be sharpened on the following interval.

\begin{theorem}[Central fixed-density profile]\label{thm:main}
Let $m\ge3$ be odd, and let $W$ be a graphon of edge density $p$.  Suppose
\begin{equation}
   \frac{5-\sqrt5}{10}\le p\le \frac{5+\sqrt5}{10}.
   \label{eq:p-range}
\end{equation}
Then
\begin{equation}
   t(C_{2m}^{\Alt},W)+t(C_{2m},W-p)
   \le (pq)^m.
   \label{eq:strong-main}
\end{equation}
Consequently,
\begin{equation}
   t(C_{2m}^{\Alt},W)\le(pq)^m.
   \label{eq:profile-main}
\end{equation}
Moreover, equality in \eqref{eq:profile-main} holds if and only if $W=p$ almost everywhere.
\end{theorem}

Since the constant graphon $W\equiv p$ has
\[
   t(C_{2m}^{\Alt},W)=p^mq^m=(pq)^m,
\]
Theorem~\ref{thm:main} determines the profile exactly on the interval \eqref{eq:p-range}.
Since $p(1-p)\le1/4$, the fixed-density bound \eqref{eq:profile-main} strengthens
\eqref{eq:universal-bound} throughout this interval, except that the two bounds agree at $p=1/2$.

The elementary equivalence
\begin{equation}
 \frac{5-\sqrt5}{10}\le p\le \frac{5+\sqrt5}{10}
 \quad\Longleftrightarrow\quad
 |2p-1|\le\sqrt{p(1-p)}
 \label{eq:range-equivalence}
\end{equation}
will be used repeatedly.  Indeed, after squaring, the latter inequality is
\[
   (2p-1)^2\le p(1-p),
\]
which is equivalent to
\[
   5p^2-5p+1\le0.
\]

\section{Finite-rank reduction}

We first reduce all determinant calculations to ordinary finite matrices.  The reduction also preserves the edge density exactly.

\begin{lemma}[Step-graphon reduction]\label{lem:step-reduction}
For fixed $m$, it is enough to prove Theorem~\ref{thm:universal} for symmetric step graphons.  For
fixed $p$ and $m$, it is enough to prove Theorem~\ref{thm:main} for symmetric step graphons of edge
density exactly $p$.
\end{lemma}

\begin{proof}
For $n\ge1$, partition $[0,1]$ into the $2^n$ dyadic intervals
\[
   I_{n,a}=\left[\frac{a}{2^n},\frac{a+1}{2^n}\right),
   \qquad 0\le a<2^n,
\]
with the harmless convention that the final interval contains $1$.  Let $\mathcal A_n$ denote the finite sigma-algebra generated by these intervals, and set
\[
   W_n:=\mathbb E[W\mid\mathcal A_n\otimes\mathcal A_n].
\]
Then $W_n$ is a symmetric step graphon, $0\le W_n\le1$, and conditional expectation preserves the integral, so
\[
   \int_{[0,1]^2}W_n=p
\]
for every $n$.  Moreover,
\[
   \|W_n-W\|_2\longrightarrow0.
\]

We verify continuity of the two densities occurring in \eqref{eq:strong-main}.  First, every factor in the alternating product \eqref{eq:alt-density} lies in $[0,1]$.  Telescoping the product over its $2m$ factors gives
\[
   \bigl|t(C_{2m}^{\Alt},W_n)-t(C_{2m}^{\Alt},W)\bigr|
   \le 2m\|W_n-W\|_1
   \le 2m\|W_n-W\|_2,
\]
because replacing $W_n$ by $W$ or $1-W_n$ by $1-W$ changes the relevant factor by the same quantity $W_n-W$ in absolute value.

Next define
\[
   K_n:=W_n-p,
   \qquad
   K:=W-p.
\]
Both kernels take values in $[-1,1]$, and $\|K_n-K\|_2=\|W_n-W\|_2$.  The same telescoping argument yields
\[
   |t(C_{2m},K_n)-t(C_{2m},K)|
   \le 2m\|K_n-K\|_1
   \le 2m\|W_n-W\|_2.
\]
Thus both densities in \eqref{eq:strong-main} converge.  The same argument, applied to $2W_n-1$
and $2W-1$, treats the signed density in \eqref{eq:universal-strong}.  Any inequality proved for
every $W_n$ therefore passes to $W$.
\end{proof}

Henceforth, until Section~\ref{sec:limit}, we assume that $W$ is a symmetric step graphon.  Let
\[
   [0,1]=B_1\sqcup\cdots\sqcup B_N
\]
be a finite measurable partition on whose product cells $W$ is constant, and let
\[
   E:=\operatorname{span}\{\mathbf 1_{B_1},\ldots,\mathbf 1_{B_N}\}.
\]
Every operator associated with a step kernel maps $L^2([0,1])$ into $E$ and vanishes on $E^\perp$.  Thus all nonzero action takes place on the finite-dimensional Hilbert space $E$.

Let
\[
   e:=\one.
\]
Because the underlying measure space has total mass one, $\|e\|_2=1$.  Let
\[
   P:=e\otimes e
\]
be the rank-one orthogonal projection onto the constants.  The operator of the constant kernel $1$ is exactly $P$.

For every symmetric step kernel $K$ and every $r\ge2$,
\begin{equation}
   t(C_r,K)=\Tr(T_K^r).
   \label{eq:trace-cycle}
\end{equation}
Likewise, direct multiplication around the alternating cycle gives
\begin{equation}
   t(C_{2m}^{\Alt},W)
   =\Tr\bigl((T_WT_{1-W})^m\bigr).
   \label{eq:trace-alt}
\end{equation}
All traces and determinants below are ordinary finite-dimensional traces and determinants on $E$.

\section{Centering and normalization at fixed density}

Let
\[
   s:=\sqrt{pq}
\]
and define the normalized centered operator
\begin{equation}
   Y:=\frac{T_W-pP}{s}.
   \label{eq:def-Y}
\end{equation}
Since the interval \eqref{eq:p-range} is contained in $(0,1)$, we have $s>0$.  The operator $Y$ is self-adjoint.  Moreover,
\begin{equation}
   \ip{e}{Ye}
   =\frac{\ip{e}{T_We}-p}{s}
   =0.
   \label{eq:mu1zero}
\end{equation}

The Hilbert--Schmidt norm of $Y$ satisfies the following bound.

\begin{lemma}[Normalized Hilbert--Schmidt bound]\label{lem:normalized-hs-bound}
One has
\begin{equation}
   \tau:=\Tr(Y^2)\le1.
   \label{eq:tau}
\end{equation}
Consequently every eigenvalue $\lambda$ of $Y$ satisfies $|\lambda|\le1$.
\end{lemma}

\begin{proof}
For a symmetric step kernel, the Hilbert--Schmidt norm agrees with the $L^2$ norm of the kernel.  Therefore
\begin{align*}
   \Tr(Y^2)
   &=\frac{1}{pq}\int_{[0,1]^2}(W(x,y)-p)^2\,dx\,dy\\
   &=\frac{1}{pq}\left(\int W^2-p^2\right).
\end{align*}
Since $0\le W\le1$, we have $W^2\le W$ pointwise, hence
\[
   \int W^2-p^2\le p-p^2=pq.
\]
This proves \eqref{eq:tau}.  Since $Y$ is self-adjoint,
\[
   \|Y\|_{\op}^2\le\|Y\|_{\HS}^2=\Tr(Y^2)\le1,
\]
so every eigenvalue belongs to $[-1,1]$.
\end{proof}

Define
\begin{equation}
   \alpha:=\sqrt{\frac qp},
   \qquad
   \delta:=\alpha^{-1}-\alpha
   =\frac{p-q}{\sqrt{pq}}.
   \label{eq:alpha-delta}
\end{equation}
The density assumption \eqref{eq:p-range}, by \eqref{eq:range-equivalence}, is precisely
\begin{equation}
   |\delta|\le1.
   \label{eq:delta-bound}
\end{equation}

Normalize the red and blue operators by their respective edge densities:
\begin{equation}
   R:=\frac{T_W}{p}=P+\alpha Y,
   \qquad
   B:=\frac{T_{1-W}}q=P-\alpha^{-1}Y.
   \label{eq:R-B}
\end{equation}
Indeed,
\[
   T_W=pP+sY=p(P+\alpha Y),
\]
because $p\alpha=s$, and similarly
\[
   T_{1-W}=P-T_W=qP-sY=q(P-\alpha^{-1}Y).
\]
The kernels of $R$ and $B$ are respectively $W/p$ and $(1-W)/q$, and hence are pointwise nonnegative.

Set
\begin{equation}
   L:=RB.
   \label{eq:def-L}
\end{equation}
Then
\begin{equation}
   T_WT_{1-W}=pq\,L,
\end{equation}
so \eqref{eq:trace-alt} becomes
\begin{equation}
   t(C_{2m}^{\Alt},W)=(pq)^m\Tr(L^m).
   \label{eq:alt-normalized}
\end{equation}
Likewise,
\begin{equation}
   T_{W-p}=T_W-pP=sY,
\end{equation}
so
\begin{equation}
   t(C_{2m},W-p)
   =(pq)^m\Tr(Y^{2m}).
   \label{eq:centered-cycle-normalized}
\end{equation}
Thus Theorem~\ref{thm:main} reduces to the finite-dimensional inequality
\begin{equation}
   \Tr(L^m)+\Tr(Y^{2m})\le1.
   \label{eq:normalized-inequality}
\end{equation}

\section{The rank-two determinant factorization}

We now derive the determinant identity used to evaluate the alternating-cycle density.
Since $Ye$ will appear repeatedly, write
\[
   u:=Ye.
\]
From \eqref{eq:R-B},
\begin{equation}
\begin{aligned}
   L
   &=(P+\alpha Y)(P-\alpha^{-1}Y)\\
   &=P-\alpha^{-1}PY+\alpha YP-Y^2.
\end{aligned}
   \label{eq:L-expanded}
\end{equation}
Because $P=e\otimes e$, self-adjointness of $Y$ gives
\[
   PY=e\otimes u,
   \qquad
   YP=u\otimes e.
\]

For an auxiliary variable $z$, define
\begin{equation}
   D(z):=I+zY^2,
   \qquad
   N(z):=D(z)^{-1}.
   \label{eq:D-N}
\end{equation}
For $z$ sufficiently close to $0$, $D(z)$ is invertible.  Define scalar functions
\begin{equation}
   h(z):=\ip{e}{N(z)e},
   \qquad
   k(z):=\ip{e}{YN(z)e},
   \qquad
   \ell(z):=\ip{e}{Y^2N(z)e}.
   \label{eq:h-k-ell}
\end{equation}
Since $N(z)$ is a function of $Y^2$, it commutes with $Y$.  The identity
\[
   N(z)+zY^2N(z)=I
\]
implies
\begin{equation}
   h(z)+z\ell(z)=1.
   \label{eq:h-ell}
\end{equation}

\begin{lemma}[Rank-two determinant factorization]\label{lem:rank-two-factorization}
With the notation above,
\begin{equation}
   \det(I-zL)
   =\det(I+zY^2)\bigl(1-zF(z)\bigr),
   \label{eq:det-factor}
\end{equation}
where
\begin{equation}
   F(z):=h(z)^2-\delta k(z)+zk(z)^2.
   \label{eq:def-F}
\end{equation}
\end{lemma}

\begin{proof}
From \eqref{eq:L-expanded},
\begin{align}
   I-zL
   &=I+zY^2-zP+z\alpha^{-1}PY-z\alpha YP\notag\\
   &=D(z)+ze(\alpha^{-1}u-e)^*-z\alpha u e^*.
   \label{eq:ranktwo-perturb}
\end{align}
Introduce the two-column operators
\[
   \mathcal U
   :=\begin{pmatrix}ze&-z\alpha u\end{pmatrix},
   \qquad
   \mathcal V
   :=\begin{pmatrix}\alpha^{-1}u-e&e\end{pmatrix}.
\]
Then \eqref{eq:ranktwo-perturb} is
\[
   I-zL=D(z)+\mathcal U\mathcal V^*.
\]
The matrix determinant lemma gives
\begin{equation}
   \det(D+\mathcal U\mathcal V^*)
   =\det(D)\det(I_2+\mathcal V^*D^{-1}\mathcal U).
   \label{eq:matrix-det-lemma}
\end{equation}
For completeness, \eqref{eq:matrix-det-lemma} follows from Sylvester's identity
$\det(I+AB)=\det(I+BA)$:
\[
   \det(D+\mathcal U\mathcal V^*)
   =\det(D)\det(I+D^{-1}\mathcal U\mathcal V^*)
   =\det(D)\det(I_2+\mathcal V^*D^{-1}\mathcal U).
\]

We now compute the $2\times2$ matrix explicitly.  Since $N=D^{-1}$ and $u=Ye$,
\[
   \ip{u}{Ne}=k,
   \qquad
   \ip{u}{Nu}=\ell.
\]
Therefore
\begin{equation}
   I_2+\mathcal V^*N\mathcal U
   =
   \begin{pmatrix}
      1+z(\alpha^{-1}k-h) & z(\alpha k-\ell)\\
      zh & 1-z\alpha k
   \end{pmatrix}.
   \label{eq:2by2}
\end{equation}
Its determinant is
\begin{align*}
 &\bigl(1+z(\alpha^{-1}k-h)\bigr)(1-z\alpha k)
   -z^2h(\alpha k-\ell)\\
 &\qquad
 =1+z\bigl((\alpha^{-1}-\alpha)k-h\bigr)
   +z^2(h\ell-k^2).
\end{align*}
Using \eqref{eq:h-ell},
\[
   z^2h\ell=zh(1-h),
\]
so the preceding expression becomes
\begin{align*}
   1+z\delta k-zh+zh-zh^2-z^2k^2
   &=1-z\bigl(h^2-\delta k+zk^2\bigr)\\
   &=1-zF(z).
\end{align*}
Substitution into \eqref{eq:matrix-det-lemma} proves \eqref{eq:det-factor}.
\end{proof}

\section{Spectral expansion of the auxiliary series}

Diagonalize the finite-dimensional self-adjoint operator $Y$.  Let
\[
   Yv_i=\lambda_i v_i
\]
for an orthonormal eigenbasis $(v_i)$, and define
\begin{equation}
   \omega_i:=|\ip{e}{v_i}|^2.
   \label{eq:omega}
\end{equation}
Then
\begin{equation}
   \omega_i\ge0,
   \qquad
   \sum_i\omega_i=1,
   \qquad
   \sum_i\omega_i\lambda_i=\ip{e}{Ye}=0,
   \label{eq:omega-identities}
\end{equation}
and
\begin{equation}
   \sum_i\lambda_i^2=\Tr(Y^2)=\tau\le1.
   \label{eq:spectrum-bound}
\end{equation}
The scalar resolvents are
\begin{equation}
   h(z)=\sum_i\frac{\omega_i}{1+z\lambda_i^2},
   \qquad
   k(z)=\sum_i\frac{\omega_i\lambda_i}{1+z\lambda_i^2}.
   \label{eq:h-k-spectral}
\end{equation}

For $n\ge0$ and real $x,y$, define
\begin{equation}
   c_n(x,y)
   :=
   \begin{cases}
   \dfrac{x^{2n+1}+y^{2n+1}}{x+y},&x+y\ne0,\\[2mm]
   (2n+1)x^{2n},&y=-x.
   \end{cases}
   \label{eq:def-cn}
\end{equation}
The second line is the continuous extension of the first.

\begin{lemma}[Properties of the divided differences]\label{lem:cn}
For every $n\ge0$ and every real $x,y$,
\begin{equation}
   c_n(x,y)\ge0.
   \label{eq:cn-positive}
\end{equation}
Moreover, for every $n\ge1$,
\begin{equation}
   (x^2+y^2)c_n(x,y)-c_{n+1}(x,y)
   =x^2y^2c_{n-1}(x,y).
   \label{eq:cn-recurrence}
\end{equation}
Finally,
\begin{equation}
   c_n(x,x)=x^{2n}.
   \label{eq:cn-diagonal}
\end{equation}
\end{lemma}

\begin{proof}
The map $t\mapsto t^{2n+1}$ is increasing on $\R$.  If $x+y\ne0$, then
\[
   c_n(x,y)
   =\frac{x^{2n+1}-(-y)^{2n+1}}{x-(-y)}
\]
is a divided difference of this increasing function, and is therefore nonnegative.  When $y=-x$, its limiting value is the derivative
\[
   (2n+1)x^{2n}\ge0.
\]
This proves \eqref{eq:cn-positive}.

For $x+y\ne0$,
\begin{align*}
 &(x^2+y^2)c_n(x,y)-c_{n+1}(x,y)\\
 &\quad=
 \frac{(x^2+y^2)(x^{2n+1}+y^{2n+1})-(x^{2n+3}+y^{2n+3})}{x+y}\\
 &\quad=
 x^2y^2\frac{x^{2n-1}+y^{2n-1}}{x+y}\\
 &\quad=x^2y^2c_{n-1}(x,y).
\end{align*}
The identity extends to $x+y=0$ by continuity.  Formula \eqref{eq:cn-diagonal} follows immediately from \eqref{eq:def-cn}.
\end{proof}

The following partial-fraction identity connects the resolvent with these divided differences:
\begin{equation}
   \frac{1+zxy}{(1+zx^2)(1+zy^2)}
   =\frac{x}{x+y}\frac1{1+zx^2}
    +\frac{y}{x+y}\frac1{1+zy^2}.
   \label{eq:partial-fraction}
\end{equation}
For $x+y=0$ the identity is understood by continuity.  Expanding the two geometric series gives
\begin{equation}
   \frac{1+zxy}{(1+zx^2)(1+zy^2)}
   =\sum_{n\ge0}(-1)^n c_n(x,y)z^n.
   \label{eq:cn-generating}
\end{equation}

We next expand the auxiliary function $F$.

\begin{lemma}[Coefficients of $F$]\label{lem:beta-expansion}
There exist real numbers $\beta_n$ such that
\begin{equation}
   F(z)=\sum_{n\ge0}(-1)^n\beta_n z^n,
   \label{eq:F-beta}
\end{equation}
where
\begin{equation}
   \beta_n
   =\sum_{i,j}\omega_i\omega_j
     c_n(\lambda_i,\lambda_j)
     \left(1-\frac{\delta}{2}(\lambda_i+\lambda_j)\right).
   \label{eq:beta-formula}
\end{equation}
\end{lemma}

\begin{proof}
Using \eqref{eq:h-k-spectral},
\begin{align*}
 h(z)^2+zk(z)^2
 &=\sum_{i,j}\omega_i\omega_j
   \frac{1+z\lambda_i\lambda_j}
   {(1+z\lambda_i^2)(1+z\lambda_j^2)}.
\end{align*}
Also, since $\sum_j\omega_j=1$,
\begin{align*}
 -\delta k(z)
 &=-\frac{\delta}{2}\sum_{i,j}\omega_i\omega_j
 \left(
   \frac{\lambda_i}{1+z\lambda_i^2}
   +\frac{\lambda_j}{1+z\lambda_j^2}
 \right)\\
 &=-\frac{\delta}{2}\sum_{i,j}\omega_i\omega_j
 \frac{(\lambda_i+\lambda_j)(1+z\lambda_i\lambda_j)}
 {(1+z\lambda_i^2)(1+z\lambda_j^2)}.
\end{align*}
Therefore
\begin{align}
 F(z)
 =\sum_{i,j}\omega_i\omega_j
 \left(1-\frac{\delta}{2}(\lambda_i+\lambda_j)\right)
 \frac{1+z\lambda_i\lambda_j}
 {(1+z\lambda_i^2)(1+z\lambda_j^2)}.
 \label{eq:F-double-sum}
\end{align}
Substituting \eqref{eq:cn-generating} into \eqref{eq:F-double-sum} and comparing coefficients proves \eqref{eq:F-beta}--\eqref{eq:beta-formula}.
\end{proof}

In the density interval \eqref{eq:p-range}, these coefficients form a decreasing nonnegative sequence.

\begin{proposition}[Monotonicity of the coefficients]\label{prop:beta-monotone}
Under the assumptions of Theorem~\ref{thm:main},
\begin{equation}
   1=\beta_0\ge\beta_1\ge\beta_2\ge\cdots\ge0.
   \label{eq:beta-monotone}
\end{equation}
More precisely,
\begin{equation}
   \beta_{n+1}\le\tau\beta_n
   \qquad(n\ge1),
   \label{eq:beta-recursive-bound}
\end{equation}
where $\tau=\Tr(Y^2)\le1$.
\end{proposition}

\begin{proof}
We divide the proof into three steps.

\medskip
\noindent\textbf{Step 1: nonnegativity and $\beta_0=1$.}
By Lemma~\ref{lem:normalized-hs-bound} and \eqref{eq:delta-bound},
\[
   |\lambda_i|\le1,
   \qquad
   |\delta|\le1.
\]
Hence
\[
   \delta\lambda_i\le |\delta|\,|\lambda_i|\le1
\]
for every $i$, and therefore
\begin{equation}
   1-\frac{\delta}{2}(\lambda_i+\lambda_j)
   =\frac{(1-\delta\lambda_i)+(1-\delta\lambda_j)}{2}
   \ge0.
   \label{eq:weight-positive}
\end{equation}
Together with $c_n\ge0$ from Lemma~\ref{lem:cn}, formula \eqref{eq:beta-formula} gives
\[
   \beta_n\ge0
\]
for every $n$.

Since $c_0(x,y)=1$, \eqref{eq:omega-identities} yields
\begin{align*}
   \beta_0
   &=\sum_{i,j}\omega_i\omega_j
      \left(1-\frac{\delta}{2}(\lambda_i+\lambda_j)\right)\\
   &=1-\delta\sum_i\omega_i\lambda_i\\
   &=1.
\end{align*}

\medskip
\noindent\textbf{Step 2: $\beta_1\le1$.}
For $r\ge0$, define
\[
   \mu_r:=\ip{e}{Y^r e}=\sum_i\omega_i\lambda_i^r.
\]
We already know $\mu_1=0$.  Since
\[
   c_1(x,y)=x^2-xy+y^2,
\]
formula \eqref{eq:beta-formula} gives
\begin{equation}
   \beta_1=2\mu_2-\delta\mu_3.
   \label{eq:beta1}
\end{equation}
Indeed, the unweighted average of $c_1$ is $2\mu_2-\mu_1^2=2\mu_2$, while
\[
   c_1(x,y)(x+y)=x^3+y^3,
\]
so the $\delta$-term contributes $-\delta\mu_3$.

We now use the pointwise nonnegativity of the red and blue kernels.  From \eqref{eq:R-B}, a direct expansion using $Pe=e$, $\ip{e}{Ye}=0$, and $Y=Y^*$ gives
\begin{equation}
   \ip{e}{RBR e}=1-2\mu_2-\alpha\mu_3,
   \label{eq:RBR}
\end{equation}
and
\begin{equation}
   \ip{e}{BRB e}=1-2\mu_2+\alpha^{-1}\mu_3.
   \label{eq:BRB}
\end{equation}
For example,
\[
   Re=e+\alpha Ye,
\]
and therefore
\[
   \ip{e}{RBR e}=\ip{Re}{B(Re)}.
\]
Expansion of the right side gives \eqref{eq:RBR}; the second identity is analogous.

Because $R$ and $B$ have pointwise nonnegative kernels, the compositions $RBR$ and $BRB$ also have pointwise nonnegative kernels.  Hence
\[
   \ip{e}{RBR e}\ge0,
   \qquad
   \ip{e}{BRB e}\ge0.
\]
Taking the convex combination with weights $q$ and $p$, and using
\[
   p\alpha^{-1}-q\alpha
   =\frac{p-q}{\sqrt{pq}}
   =\delta,
\]
we obtain
\begin{align*}
 0
 &\le q\ip{e}{RBR e}+p\ip{e}{BRB e}\\
 &=1-2\mu_2+\delta\mu_3\\
 &=1-\beta_1.
\end{align*}
Thus
\[
   0\le\beta_1\le1.
\]

\medskip
\noindent\textbf{Step 3: the recursive inequality for $n\ge1$.}
Fix indices $i,j$.  If $i=j$, then by \eqref{eq:cn-diagonal},
\[
   c_{n+1}(\lambda_i,\lambda_i)
   =\lambda_i^2c_n(\lambda_i,\lambda_i)
   \le\tau c_n(\lambda_i,\lambda_i),
\]
because $\lambda_i^2\le\sum_r\lambda_r^2=\tau$.

If $i\ne j$, then Lemma~\ref{lem:cn} gives
\begin{align*}
   c_{n+1}(\lambda_i,\lambda_j)
   &=(\lambda_i^2+\lambda_j^2)c_n(\lambda_i,\lambda_j)
     -\lambda_i^2\lambda_j^2c_{n-1}(\lambda_i,\lambda_j)\\
   &\le(\lambda_i^2+\lambda_j^2)c_n(\lambda_i,\lambda_j)\\
   &\le\tau c_n(\lambda_i,\lambda_j),
\end{align*}
because $c_{n-1}\ge0$ and, for distinct basis indices,
\[
   \lambda_i^2+\lambda_j^2\le\sum_r\lambda_r^2=\tau.
\]

The factors \eqref{eq:weight-positive} in \eqref{eq:beta-formula} are nonnegative.  Multiplying the preceding pointwise inequality by those factors and by $\omega_i\omega_j$, then summing over $i,j$, yields
\[
   \beta_{n+1}\le\tau\beta_n.
\]
Since $\tau\le1$, this proves
\[
   \beta_{n+1}\le\beta_n
\]
for every $n\ge1$.  Together with $1=\beta_0\ge\beta_1\ge0$, the full chain \eqref{eq:beta-monotone} follows.
\end{proof}

\section{The odd logarithmic coefficient lemma}

The preceding proposition reduces the remaining argument to a one-variable formal power-series fact.

\begin{lemma}[Odd logarithmic coefficient]\label{lem:odd-log}
Let
\[
   1=\beta_0\ge\beta_1\ge\beta_2\ge\cdots\ge0
\]
and define
\[
   F(z)=\sum_{n\ge0}(-1)^n\beta_nz^n.
\]
Then for every odd positive integer $m$,
\begin{equation}
   m[z^m]\bigl(-\log(1-zF(z))\bigr)\le1.
   \label{eq:odd-log-bound}
\end{equation}
More precisely, if
\[
   d_n:=\beta_n-\beta_{n+1}\ge0,
   \qquad
   G(z):=\sum_{n\ge0}d_nz^n,
\]
then
\begin{equation}
\begin{aligned}
 1-m[z^m]\bigl(-\log(1-zF(z))\bigr)
 =m\sum_{r=1}^{(m-1)/2}\frac1r
   [z^{m-2r}]G(z)^r
 \ge0.
   \label{eq:odd-log-identity}
\end{aligned}
\end{equation}
\end{lemma}

\begin{proof}
For $k\ge1$, the coefficient of $z^k$ in $(1+z)F(z)$ is
\[
   (-1)^k\beta_k+(-1)^{k-1}\beta_{k-1}
   =(-1)^{k-1}d_{k-1}.
\]
Thus
\begin{equation}
   (1+z)F(z)=1+zG(-z).
   \label{eq:FG}
\end{equation}
It follows that
\begin{align}
   (1+z)(1-zF(z))
   &=1+z-z(1+z)F(z)\notag\\
   &=1-z^2G(-z).
   \label{eq:factor-log}
\end{align}
Taking formal logarithms gives
\begin{equation}
   -\log(1-zF(z))
   =\log(1+z)-\log(1-z^2G(-z)).
   \label{eq:log-split}
\end{equation}
Since $m$ is odd,
\[
   [z^m]\log(1+z)=\frac1m.
\]
Also,
\begin{equation}
   -\log(1-z^2G(-z))
   =\sum_{r\ge1}\frac{z^{2r}G(-z)^r}{r}.
   \label{eq:second-log}
\end{equation}
For $1\le r\le(m-1)/2$, the integer $m-2r$ is odd.  Since $G$ has nonnegative coefficients,
\[
   [z^{m-2r}]G(-z)^r
   =(-1)^{m-2r}[z^{m-2r}]G(z)^r
   =-[z^{m-2r}]G(z)^r.
\]
Therefore the coefficient of $z^m$ in the second term on the right side of \eqref{eq:log-split} is
\[
   -\sum_{r=1}^{(m-1)/2}\frac1r[z^{m-2r}]G(z)^r.
\]
Combining this with the contribution $1/m$ from $\log(1+z)$ proves the exact identity \eqref{eq:odd-log-identity}, and hence \eqref{eq:odd-log-bound}.
\end{proof}

\section{Completion of the finite-dimensional argument}

We now combine the determinant factorization with Lemma~\ref{lem:odd-log}.  For any finite matrix $M$, the Taylor expansion at $z=0$ gives
\begin{equation}
   -\log\det(I-zM)
   =\sum_{r\ge1}\frac{\Tr(M^r)}{r}z^r.
   \label{eq:trace-log}
\end{equation}
This identity may equivalently be read as a formal power-series identity.  Applying \eqref{eq:trace-log} to Lemma~\ref{lem:rank-two-factorization},
\begin{equation}
   -\log\det(I-zL)
   =-\log\det(I+zY^2)-\log(1-zF(z)).
   \label{eq:log-det-factor}
\end{equation}
The first term on the right has expansion
\begin{equation}
   -\log\det(I+zY^2)
   =\sum_{r\ge1}\frac{(-1)^r\Tr(Y^{2r})}{r}z^r.
   \label{eq:Y-log}
\end{equation}
Taking the coefficient of $z^m$ in \eqref{eq:log-det-factor}, and using that $m$ is odd, gives
\begin{equation}
   \Tr(L^m)
   =-\Tr(Y^{2m})
    +m[z^m]\bigl(-\log(1-zF(z))\bigr).
   \label{eq:trace-identity-final}
\end{equation}
By Proposition~\ref{prop:beta-monotone}, the coefficients of $F$ satisfy the hypotheses of Lemma~\ref{lem:odd-log}.  Hence
\[
   m[z^m]\bigl(-\log(1-zF(z))\bigr)\le1.
\]
Substitution into \eqref{eq:trace-identity-final} yields
\begin{equation}
   \Tr(L^m)+\Tr(Y^{2m})\le1.
   \label{eq:abstract-proved}
\end{equation}
This proves the normalized inequality \eqref{eq:normalized-inequality}.  Multiplying by $(pq)^m$ and using \eqref{eq:alt-normalized} and \eqref{eq:centered-cycle-normalized}, we obtain
\[
   t(C_{2m}^{\Alt},W)+t(C_{2m},W-p)\le(pq)^m
\]
for every symmetric step graphon of edge density $p$ satisfying \eqref{eq:p-range}.

\section{The universal inequality}\label{sec:universal-proof}

We now prove the density-independent inequality \eqref{eq:universal-strong}, using the same
rank-two determinant factorization and scalar coefficient argument.  By Lemma~\ref{lem:step-reduction}, it is enough for
now to take $W$ to be a symmetric step graphon.  On the finite-dimensional space $E$ from
Section~2, define
\[
   X:=T_{2W-1},
   \qquad
   L_0:=(P+X)(P-X).
\]
Then $X=X^*$ and
\begin{equation}
   T_W=\frac{P+X}{2},
   \qquad
   T_{1-W}=\frac{P-X}{2}.
   \label{eq:universal-coordinates}
\end{equation}
Consequently,
\begin{equation}
   4^m t(C_{2m}^{\Alt},W)=\Tr(L_0^m),
   \qquad
   t(C_{2m},2W-1)=\Tr(X^{2m}).
   \label{eq:universal-traces}
\end{equation}
Moreover,
\begin{equation}
   \tau_0:=\Tr(X^2)
   =\int_{[0,1]^2}(2W(x,y)-1)^2\,dx\,dy
   \le1.
   \label{eq:universal-hs-bound}
\end{equation}

Set
\[
   D_0(z):=I+zX^2,
   \qquad N_0(z):=D_0(z)^{-1},
\]
and define
\[
   h_0(z):=\ip{e}{N_0(z)e},
   \qquad
   k_0(z):=\ip{e}{XN_0(z)e}.
\]
The proof of Lemma~\ref{lem:rank-two-factorization}, with $Y=X$ and $\alpha=1$ (and hence $\delta=0$),
gives
\begin{equation}
   \det(I-zL_0)
   =\det(I+zX^2)\bigl(1-zF_0(z)\bigr),
   \qquad
   F_0(z):=h_0(z)^2+zk_0(z)^2.
   \label{eq:universal-det-factor}
\end{equation}

Diagonalize $X$, writing $Xv_i=\lambda_i v_i$ and
$\omega_i:=|\ip{e}{v_i}|^2$.  Thus
\[
   \omega_i\ge0,
   \qquad
   \sum_i\omega_i=1,
   \qquad
   \sum_i\lambda_i^2=\tau_0\le1.
\]
The calculation of Lemma~\ref{lem:beta-expansion}, now with $\delta=0$, gives
\begin{equation}
   F_0(z)=\sum_{n\ge0}(-1)^n\beta_n^{(0)}z^n,
   \qquad
   \beta_n^{(0)}
   :=\sum_{i,j}\omega_i\omega_jc_n(\lambda_i,\lambda_j).
   \label{eq:universal-beta}
\end{equation}
Lemma~\ref{lem:cn} implies $\beta_n^{(0)}\ge0$, and $c_0\equiv1$ gives
$\beta_0^{(0)}=1$.

It remains to check the first monotonicity step.  Define
\[
   \mu_r:=\sum_i\omega_i\lambda_i^r=\ip{e}{X^re}.
\]
Since $c_1(x,y)=x^2-xy+y^2$,
\begin{equation}
   \beta_1^{(0)}=2\mu_2-\mu_1^2.
   \label{eq:universal-beta-one}
\end{equation}
Decompose $E=\R e\oplus e^\perp$ and write
\[
   X=
   \begin{pmatrix}
      a&w^*\\
      w&A
   \end{pmatrix}.
\]
Then $\mu_1=a$, $\mu_2=a^2+\|w\|^2$, and hence
\[
   \beta_1^{(0)}=a^2+2\|w\|^2
   \le a^2+2\|w\|^2+\Tr(A^2)
   =\tau_0
   \le1.
\]
For $n\ge1$, Step~3 in the proof of Proposition~\ref{prop:beta-monotone} applies unchanged, with
the additional density factor in \eqref{eq:beta-formula} equal to one, and gives
\[
   \beta_{n+1}^{(0)}\le\tau_0\beta_n^{(0)}\le\beta_n^{(0)}.
\]
Therefore
\[
   1=\beta_0^{(0)}\ge\beta_1^{(0)}\ge\beta_2^{(0)}\ge\cdots\ge0.
\]

Lemma~\ref{lem:odd-log}, applied to \eqref{eq:universal-beta}, yields
\[
   m[z^m]\bigl(-\log(1-zF_0(z))\bigr)\le1.
\]
Finally, comparing the coefficient of $z^m$ in the logarithm of
\eqref{eq:universal-det-factor}, exactly as in Section~7, gives
\[
   \Tr(L_0^m)+\Tr(X^{2m})
   =m[z^m]\bigl(-\log(1-zF_0(z))\bigr)
   \le1.
\]
Together with \eqref{eq:universal-traces}, this proves \eqref{eq:universal-strong} for symmetric
step graphons.

\section{Passage to arbitrary graphons and equality}\label{sec:limit}

By Lemma~\ref{lem:step-reduction}, the inequalities \eqref{eq:universal-strong} and
\eqref{eq:strong-main} pass from step graphons to arbitrary graphons.

We use the following standard positivity fact.  For every bounded symmetric kernel $K$ and every
integer $r\ge1$,
\begin{equation}
   t(C_{2r},K)=\Tr(T_K^{2r})=\|T_K^r\|_{\HS}^2\ge0.
   \label{eq:even-signed-nonnegative}
\end{equation}
For step kernels this follows from
$T_K^{2r}=(T_K^r)^*T_K^r$ and the finite-dimensional trace formula.  For a general bounded
symmetric kernel, apply the dyadic conditional expectations from Lemma~\ref{lem:step-reduction}.
The cycle densities converge by the same telescoping argument, while the corresponding integral
operators and their $r$th powers converge in Hilbert--Schmidt norm.  Passing to the limit gives
\eqref{eq:even-signed-nonnegative}.

Applying \eqref{eq:even-signed-nonnegative} to $K=2W-1$ in
\eqref{eq:universal-strong} gives \eqref{eq:universal-bound}.  If equality holds, then
\eqref{eq:universal-strong} forces
\[
   \|T_{2W-1}^m\|_{\HS}^2=t(C_{2m},2W-1)=0.
\]
Because $T_{2W-1}$ is self-adjoint, this implies $T_{2W-1}=0$.  Its Hilbert--Schmidt norm is the
$L^2$ norm of its kernel, so $W=1/2$ almost everywhere.  The converse is immediate.  This proves
Theorem~\ref{thm:universal}.

Applying \eqref{eq:even-signed-nonnegative} to $K=W-p$ in \eqref{eq:strong-main} gives
\[
   t(C_{2m}^{\Alt},W)\le(pq)^m,
\]
which is \eqref{eq:profile-main}.

Suppose equality holds in \eqref{eq:profile-main}.  Then \eqref{eq:strong-main} and \eqref{eq:even-signed-nonnegative} force
\[
   t(C_{2m},W-p)=0.
\]
Hence every eigenvalue of the self-adjoint compact operator $T_{W-p}$ is zero.  Therefore $T_{W-p}=0$.  Since $T_{W-p}$ is Hilbert--Schmidt,
\[
   \|W-p\|_2^2=\|T_{W-p}\|_{\HS}^2=0,
\]
so
\[
   W=p
\]
almost everywhere.  Conversely, the constant graphon $W\equiv p$ plainly gives equality.

We have therefore also proved Theorem~\ref{thm:main} in full.

\begin{corollary}[Exact profile on the central interval]
Let $m\ge3$ be odd.  For
\[
   \frac{5-\sqrt5}{10}\le p\le\frac{5+\sqrt5}{10},
\]
one has
\[
   \sup\left\{
      t(C_{2m}^{\Alt},W):
      W\text{ a graphon and }t(K_2,W)=p
   \right\}
   =[p(1-p)]^m.
\]
The unique maximizing graphon, up to equality almost everywhere, is the constant graphon $W\equiv p$.
\end{corollary}

\section{Remark on the Lean verification}

The inequalities \eqref{eq:universal-strong}, \eqref{eq:universal-bound},
\eqref{eq:strong-main}, and \eqref{eq:profile-main}, their cyclic-integral formulations, and
attainment by the constant graphons $W\equiv1/2$ and $W\equiv p$ have been formalised in the
Lean~4 proof assistant~\cite{Lean4}, on top of the mathlib library~\cite{mathlib}.  The uniqueness
of the equality cases is not formalised.  The formal proof differs from the one given here in
three ways.

First, instead of the step-graphon approximation (Lemma~\ref{lem:step-reduction}), we write the
graphon quantities in \eqref{eq:universal-strong} and \eqref{eq:strong-main}, the latter after
division by $(pq)^m$, and their matrix counterparts as the same polynomials in the moments
$\mu_g=\ip{\one}{Z^g\one}$, $g<2m$, where $Z$ denotes the relevant normalized kernel operator.
We then compress $Z$ to the Krylov subspace
$\operatorname{span}\{\one,Z\one,\ldots,Z^{2m}\one\}$; this produces a finite symmetric matrix
with the same moments and $\Tr(A^2)\le1$, to which the corresponding finite-dimensional inequality
applies directly.

Second, the formal proof never forms a logarithm or a large determinant.  Instead of handling
$m[z^m]\bigl(-\log A(z)\bigr)$, it uses the operator
$\Lambda(A):=-zA'(z)/A(z)$: since $z\,d/dz$ multiplies the coefficient of $z^m$ by $m$, one has
$[z^m]\Lambda(A)=m[z^m]\bigl(-\log A(z)\bigr)$, while $\Lambda$ itself involves no logarithm.  The
rule $\Lambda(fg)=\Lambda(f)+\Lambda(g)$, immediate from $(fg)'=f'g+fg'$, replaces the only
property of the logarithm that the argument uses.  Similarly, the logarithm in
\eqref{eq:trace-identity-final} is read as $[z^m]\Lambda(1-zF(z))$, and the identity is proved not
through \eqref{eq:trace-log} but directly from the rank-two structure in the proof of
Lemma~\ref{lem:rank-two-factorization}, as a computation with the $2\times2$ matrix
$I_2+\mathcal V^*N(z)\mathcal U$; the only determinants that appear are $2\times2$.

Third, for the first coefficient in the universal inequality, the formal proof obtains
$\beta_1^{(0)}\le\tau_0$ from a direct sum-of-squares identity, without decomposing $X$ along $e$.

\begin{thebibliography}{9}

\bibitem{BGHKS}
A. Basit, B. Granet, D. Horsley, A. K\"undgen, and K. Staden,
\emph{The semi-inducibility problem},
arXiv:2501.09842.

\bibitem{ChenNoel}
H. Chen and J. A. Noel,
\emph{On alternating 6-cycles in edge-coloured graphs},
arXiv:2505.09809.

\bibitem{Lovasz}
L. Lov\'asz,
\emph{Large Networks and Graph Limits},
American Mathematical Society Colloquium Publications, vol.~60,
American Mathematical Society, Providence, RI, 2012.

\bibitem{mathlib}
The mathlib Community,
\emph{The Lean mathematical library},
Proceedings of the 9th ACM SIGPLAN International Conference on
Certified Programs and Proofs (CPP 2020), ACM, 2020, pp.~367--381.

\bibitem{Lean4}
L. de Moura and S. Ullrich,
\emph{The Lean 4 theorem prover and programming language},
Automated Deduction -- CADE 28, Lecture Notes in Computer Science,
vol.~12699, Springer, 2021, pp.~625--635.

\end{thebibliography}

\end{document}
